\documentclass[11pt,a4paper]{article}
\usepackage[T1]{fontenc}
\usepackage[utf8]{inputenc}
\usepackage{lmodern}
\usepackage[margin=1.9cm]{geometry}
\usepackage{booktabs,enumitem,xcolor,titlesec}
\usepackage[colorlinks=true,urlcolor=blue!50!black]{hyperref}
\pagestyle{empty}
\setlength{\parindent}{0pt}\setlength{\parskip}{0.35em}
\setlist{nosep,leftmargin=1.1em}
\titleformat{\section}{\large\bfseries\color{blue!45!black}}{}{0pt}{}
\titlespacing*{\section}{0pt}{0.9em}{0.3em}

\begin{document}
{\LARGE\bfseries Cheat sheet}\hfill{\small Chest X-ray Diagnosis}\par
{\small Read once. Everything here is true and checked.}

\section*{The pitch (30 seconds)}
\begin{itemize}
  \item We detect \textbf{14 chest diseases} from an X-ray, several can be present at once.
  \item A DenseNet121 (pretrained on ImageNet, fine-tuned) gets \textbf{0.799 AUC} on the official test set.
  \item Grad-CAM heatmaps show where the model looked. It runs as a web app with tests and CI.
\end{itemize}

\section*{Numbers to remember}
\begin{tabular}{@{}ll@{\qquad}ll@{}}
Images & 112,120 (30,805 patients) & Val AUC & \textbf{0.831} \\
Train / val / test & 73,916 / 12,608 / 25,596 & Test AUC & \textbf{0.799} \\
Model size & 7.0\,M parameters & Best / worst class & Hernia 0.93 / Pneumonia 0.69 \\
Training & 15 epochs, Adam, lr $10^{-4}$ & Speed & 25\,ms per image (CPU) \\
\end{tabular}

\section*{What is ours, honestly}
\begin{itemize}
  \item The \textbf{model was trained by the team} in the notebook (15 epochs). We did not invent a new architecture.
  \item We added: a reproducible, tested pipeline, honest test evaluation, Grad-CAM, ONNX, the API, the UI, Docker, CI, DVC, and the report.
  \item The \textbf{ablations} are new training runs (5 epochs each) on Kaggle.
\end{itemize}

\section*{Questions you will get}
\textbf{Why split by patient?} One patient has up to 184 images. A random split puts the same person in train and validation, so the model recognises the person instead of the disease and the score is inflated.

\textbf{Why AUC and not accuracy?} Rare diseases (Hernia 0.15\%): predicting ``healthy'' always is 99\% accurate and useless. AUC measures how well the model ranks sick above healthy.

\textbf{Why weighted loss?} The idea: rare diseases would be ignored, so each positive counts $N_{neg}/N_{pos}$ times (5 to 653). \textbf{But our ablation showed it did not help AUC} (plain loss 0.836 vs weighted 0.823, better in 13 of 14 classes) and it makes scores over-confident. Say this openly: it is a real finding.

\textbf{What is an ablation?} Retrain with one choice removed and compare, to measure what that choice is worth.

\textbf{Overfitting?} Test gap is small (0.831 vs 0.799, best epoch picked on validation). In the 5-epoch reference run, validation loss rises after epoch 2--3 while training loss keeps falling, and AUC peaks at epoch 4: overfitting starts early. We have no curve for the 15-epoch run.

\textbf{Is 0.799 good?} Reasonable for plain DenseNet121 at 224\,px. Published results are a bit higher. Next steps: higher resolution, stronger backbone.

\textbf{Why is Pneumonia weak?} Diffuse and ambiguous on a frontal film; labels were machine-extracted from reports (noisy).

\textbf{Usable clinically?} No. Noisy labels, one dataset, no external validation, one run, no calibration.

\section*{Ablation results (5 epochs each, validation AUC)}
\begin{tabular}{@{}lr@{\quad}l@{}}
Fine-tuned, weighted loss (reference) & 0.823 & our recipe \\
Frozen backbone, last layer only & 0.734 & fine-tuning is essential (+0.09) \\
Fine-tuned, plain loss & \textbf{0.836} & better than our weighted loss \\
\end{tabular}
\begin{itemize}
  \item \textbf{Fine-tuning matters a lot.} Frozen still improves but is far behind.
  \item \textbf{Weighting did not help} (+0.013 for plain loss, every epoch, 13/14 classes). One seed, so suggestive, not proven. A 5-epoch plain-loss run already beats the 15-epoch main model's validation AUC (0.836 vs 0.831); not tested on the test set.
  \item Strong augmentation was configured but dropped to save GPU time: no claim about it.
\end{itemize}

\section*{Honest weak spots (we measured them)}
\begin{itemize}
  \item \textbf{Precision is low.} At a cutoff that catches 82\% of sick cases (mean sensitivity), specificity is 0.62 and precision only 0.14 on average. Good at ranking, noisy as a yes/no alarm.
  \item \textbf{Scores are not probabilities.} Mean predicted 0.38 vs 0.076 real positives: the weighted loss makes the model over-confident. The app's ``99.9\%'' means ``scores very high'', not ``99.9\% sure''.
  \item Hernia has the best AUC (0.93) but only 0.3\% prevalence in test, so precision is tiny.
\end{itemize}

\section*{Do not say}
\begin{tabular}{@{}p{0.46\textwidth}p{0.48\textwidth}@{}}
\toprule
Don't say & Say instead \\
\midrule
``It diagnoses diseases'' & ``Research demo; it ranks likelihoods, not a diagnosis'' \\
``99.9\% sure'' & ``It scores very high; scores are not calibrated'' \\
``Grad-CAM proves it looks at the right place'' & ``Plausibility check; localisation was not measured'' \\
``We built the model'' & ``The team trained it; we built the pipeline around it'' \\
``The webcam works'' & ``Implemented; not yet tested on a real camera'' (until you test it) \\
\bottomrule
\end{tabular}

\section*{Demo (2 minutes)}
Open the app $\to$ click the Cardiomegaly example $\to$ read bars and latency $\to$ click Cardiomegaly for the heatmap $\to$ show Hernia is low (it does not flag everything). Say the gallery images are confident test cases.

\vfill
{\footnotesize Repo: github.com/khldd/chest-xray-diagnosis \quad|\quad Full report: report/main.pdf \quad|\quad Visuals: report/visuals.pdf}
\end{document}
